\documentclass[letterpaper,11pt]{article}

%-------------------------
% PACKAGES
%-------------------------
\usepackage{latexsym}
\usepackage[empty]{fullpage}
\usepackage{titlesec}
\usepackage{marvosym}
\usepackage[usenames,dvipsnames]{color}
\usepackage{verbatim}
\usepackage{enumitem}
\usepackage[pdftex]{hyperref}
\usepackage{fancyhdr}
\usepackage{graphicx}
\usepackage[left=0.8cm,right=0.8cm,top=0.5cm,bottom=1.5cm]{geometry}

%-------------------------
% COLORS
%-------------------------
\definecolor{softgray}{RGB}{95,95,95}
\definecolor{lightgray}{RGB}{130,130,130}

%-------------------------
% PAGE SETUP
%-------------------------
\pagestyle{fancy}
\fancyhf{}
\fancyfoot{}
\renewcommand{\headrulewidth}{0pt}
\renewcommand{\footrulewidth}{0pt}

\urlstyle{same}
\raggedbottom
\raggedright
\setlength{\tabcolsep}{0in}

%-------------------------
% SECTION FORMAT
%-------------------------
\titleformat{\section}
{\scshape\raggedright\large}
{}{0em}{}
[\color{lightgray}\titlerule \vspace{-4pt}]

%-------------------------
% CUSTOM COMMANDS
%-------------------------
\newcommand{\resumeItem}[1]{
  \item\small{#1}
}

\newcommand{\resumeSubheading}[4]{
  \vspace{2pt}\item
    \begin{tabular*}{0.97\textwidth}{l@{\extracolsep{\fill}}r}
      \textbf{#1} & \textit{\small #2} \\
      \textit{\small#3} & \textit{\small #4} \\
    \end{tabular*}\vspace{-3pt}
}

\newcommand{\resumeSubHeadingListStart}{\begin{itemize}[leftmargin=*, itemsep=1pt, topsep=2pt]}
\newcommand{\resumeSubHeadingListEnd}{\end{itemize}}
\newcommand{\resumeItemListStart}{\begin{itemize}[itemsep=1pt, topsep=1pt]}
\newcommand{\resumeItemListEnd}{\end{itemize}\vspace{-6pt}}
\newcommand{\resumeSubItem}[2]{
  \item\small{\textbf{#1}{: #2}}\vspace{-1pt}
}

\hypersetup{
    colorlinks=true,
    linkcolor=softgray,
    urlcolor=softgray
}

%=========================
% DOCUMENT
%=========================
\begin{document}

%----------HEADER-----------------
\begin{tabular*}{\textwidth}{@{\extracolsep{\fill}}l r}

\begin{minipage}[t]{0.6\textwidth}
\raggedright
{\LARGE\bfseries Théo Pantecouteau} \\
{\color{softgray}\small
\href{https://www.linkedin.com/in/theo-pantecouteau/}{linkedin} --
\href{https://github.com/tpantecouteau}{github}} \\
{\color{lightgray}\small Paris, France}
\end{minipage}
&
\begin{minipage}[t]{0.35\textwidth}
\raggedleft
{\small
06 51 06 25 35 \\
\href{mailto:theopantecouteau1@protonmail.com}{theopantecouteau1@protonmail.com}}
\end{minipage}

\end{tabular*}

\vspace{-4mm}

%-----------PROFILE-----------------
\section{Profil}
\small{
Data Engineer avec forte expertise en développement logiciel, spécialisé en \textbf{pipelines de données}, \textbf{APIs backend} et \textbf{intégration de solutions IA}. Expérience en architecture ETL, automatisation de workflows analytiques.
}
\vspace{-4mm}

%-----------EXPERIENCE-----------------
\section{Expériences professionnelles}
\resumeSubHeadingListStart

\resumeSubheading
{Data Engineer Python — Pipelines \& Automatisation}{Oct. 2023 -- Oct. 2025}
{Ministère de l'Intérieur}{Paris, France}
\resumeItemListStart
\resumeItem{Conception et production de \textbf{pipelines ETL Python/SQL} sous Dataiku servant 30+ utilisateurs métier quotidiennement, avec automatisation complète du workflow analytique (réduction temps de traitement : 3 jours → 30 min, élimination interventions manuelles).}
\resumeItem{Architecture d'un système de \textbf{monitoring et quality checks} : logs structurés, contrôles de cohérence automatisés (mensuels/annuels), gestion d'erreurs et alerting pour garantir la fiabilité des données en production.}
\resumeItem{Analyses exploratoires et aide à la décision : clustering k-means, segmentation de données et dashboards automatisés pour insights métier.}
\resumeItem{Déploiement en environnement institutionnel contraint avec exigences élevées de traçabilité, conformité et continuité de service.}
\resumeItemListEnd

\vspace{2mm}

\resumeSubheading
{Software Engineer Python — Backend \& Data Ingestion}{Sep. 2022 -- Oct. 2023}
{Cash Flow Positif}{Paris, France}
\resumeItemListStart
\resumeItem{Développement d'\textbf{APIs REST robustes} (FastAPI) avec pipeline d'ingestion multi-sources : scraping automatisé, traitement de 4\,000+ annonces/jour depuis ~10 sources hétérogènes, gestion d'erreurs et reprise sur échec.}
\resumeItem{Modélisation de \textbf{schémas PostgreSQL} pour données immobilières : structuration, normalisation et optimisation pour requêtes analytiques et applicatives.}
\resumeItem{Stack fullstack : automatisation de workflows métier (publication annonces, messaging, scheduling) avec intégration continue.}
\resumeItemListEnd

\resumeSubHeadingListEnd
\vspace{-5mm}

%-----------PROJECTS-----------------
\section{Projets Data \& IA}
\resumeSubHeadingListStart

\resumeSubItem{Plateforme RAG pour publications scientifiques}
{Développement d'une solution \textbf{Retrieval-Augmented Generation} pour exploration de corpus ArXiv. Pipeline complet : ingestion PDF, chunking optimisé, \textbf{embeddings sémantiques}, indexation vectorielle (Pinecone), requêtes contextualisées via LLM. Permet interrogation en langage naturel de documents techniques complexes. \href{https://github.com/tpantecouteau/arxiv_cs_paper_rag}{\underline{Code source}}.}

\resumeSubItem{ReadLater AI — Application fullstack avec infra production}
{Application d'agrégation de contenus web avec résumés IA. \textbf{Stack} : FastAPI, Next.js 15, PostgreSQL, extension Chrome (Manifest V3). \textbf{Infrastructure} : déploiement AWS (EC2, RDS, VPC) via Terraform, Docker multi-stage builds. \textbf{Sécurité production} : reverse proxy Nginx (rate limiting, security headers), containers hardened (read-only, non-root, no-new-privileges), Fail2ban, audit et patch de dépendances (pip-audit, pnpm audit). \href{https://github.com/tpantecouteau/readlater-ai}{\underline{Code source}}.}


\resumeSubItem{Secure Drop — Service zero-knowledge de partage de fichiers}
{Architecture backend sécurisée pour transfert de fichiers avec chiffrement bout-en-bout. Le serveur n'accède jamais aux clés ni au contenu. API REST dédiée, réflexion sur modèle de menace et garanties cryptographiques. \href{https://github.com/tpantecouteau/secure_drop}{\underline{Code source}}.}

\resumeSubHeadingListEnd
\vspace{-6mm}

%-----------EDUCATION-----------------
\section{Parcours scolaire}
\resumeSubHeadingListStart
\resumeSubheading
{ESILV}{Paris, France}
{Diplôme d'ingénieur en Informatique — Data Science \& IA}{2022 -- 2025}
\resumeSubheading
{IUT Montpellier-Sète}{Montpellier, France}
{DUT Informatique}{2020 -- 2022}
\resumeSubHeadingListEnd
\vspace{-5mm}

%-----------SKILLS-----------------
\section{Compétences}
\resumeSubHeadingListStart
\item{\textbf{Data Engineering} : Pipelines ETL, Dataiku, ingestion multi-sources, scraping à grande échelle, data quality}
\item{\textbf{Backend \& APIs} : Python, FastAPI, REST, OpenAPI, architectures robustes}
\item{\textbf{AI/ML} : RAG, intégration LLMs, embeddings sémantiques, prompt engineering, vector databases (Pinecone)}
\item{\textbf{Databases} : PostgreSQL (modélisation, optimisation), vector stores, schéma design}
\item{\textbf{Cloud \& Production} : AWS (EC2, RDS, VPC, S3, Lambda), Docker, Terraform, CI/CD}
\item{\textbf{DevOps \& Sécurité} : Nginx (reverse proxy, rate limiting), container hardening, Fail2ban, audit de dépendances}
\item{\textbf{Observabilité} : Logs structurés, error handling, retries, alerting, monitoring (Prometheus, Grafana, CloudWatch)}
\item{\textbf{Langages} : Python (expert), SQL, JavaScript/TypeScript}
\item{\textbf{Langues} : Français (natif), Anglais B2+ (TOEIC 935)}
\resumeSubHeadingListEnd

\end{document}